%% ma: L12-B19-0001
%% lop: 12
%% chuong: 6
%% bai: 19
%% dang: 12.19.1
%% loai: DS
%% muc: [NB, TH, VD, VD]
%% dap_an: "ĐSĐS"
%% nguon: "0. ĐỀ THAM KHẢO TN THPT 2025.pdf (D00), Phần II Câu 3"
%% kien_thuc: "Bài 19 (công thức xác suất toàn phần và công thức Bayes), Bài 18 (xác suất có điều kiện)"
%% trang_thai: chinh-thuc
\begin{ex}
Trước khi đưa một loại sản phẩm ra thị trường, người ta đã phỏng vấn ngẫu nhiên $200$ khách hàng về sản phẩm đó. Kết quả thống kê như sau: có $105$ người trả lời ``sẽ mua''; có $95$ người trả lời ``không mua''. Kinh nghiệm cho thấy tỉ lệ khách hàng thực sự sẽ mua sản phẩm tương ứng với những cách trả lời ``sẽ mua'' và ``không mua'' lần lượt là $70\%$ và $30\%$.

Gọi $A$ là biến cố ``Người được phỏng vấn thực sự sẽ mua sản phẩm''. Gọi $B$ là biến cố ``Người được phỏng vấn trả lời sẽ mua sản phẩm''.
\choiceTF
{\True Xác suất $P(B)=\dfrac{21}{40}$ và $P(\overline B)=\dfrac{19}{40}$.}
{Xác suất có điều kiện $P(A\mid B)=0{,}3$.}
{\True Xác suất $P(A)=0{,}51$.}
{Trong số những người được phỏng vấn thực sự sẽ mua sản phẩm có $70\%$ người đã trả lời ``sẽ mua'' khi được phỏng vấn (kết quả tính theo phần trăm được làm tròn đến hàng đơn vị).}
\loigiai{a) $P(B)=\frac{105}{200}=\frac{21}{40}$, $P(\overline B)=\frac{95}{200}=\frac{19}{40}$. Đúng.
b) $P(A\mid B)=0{,}7$ và $P(A\mid\overline B)=0{,}3$. Ý b) sai.
c) $P(A)=P(B)P(A\mid B)+P(\overline B)P(A\mid\overline B)=\frac{21}{40}\cdot0{,}7+\frac{19}{40}\cdot0{,}3=0{,}51$. Đúng.
d) $P(B\mid A)=\dfrac{P(B)P(A\mid B)}{P(A)}=\dfrac{0{,}525\cdot0{,}7}{0{,}51}\approx0{,}72=72\%\ne70\%$. Sai.}
\end{ex}
